This paper studies a generalization failure discovered during an automated architecture search for NanoGPT, which we call \emph{forking}. When the training data is repeatedly consumed, models with an over-encoding n-gram memory branch exhibit a clear separation between training and validation loss, and this separation often grows in a staircase pattern aligned with epoch boundaries. We study this phenomenon in a controlled vanilla NanoGPT setting and reproduce it in a DeepSeek-style model with an Engram-like memory module. Mechanistically, we find that repeated updates sharpen the continuations observed in training and keep suppressing the probability of unseen continuations, so their loss grows pass after pass. The n-gram module partitions the model into weakly interacting context-specific subspaces, which makes this overfitting and sharpening especially pronounced. Low-frequency n-gram contexts contribute most of the gap, whereas larger training budgets and heavily crowded over-encoding tables suppress it, so forking goes unnoticed in conventional pretraining. These trends predict that short-budget, heavily repeated regimes such as SFT and RL should be prone to forking, which we verify even when the n-gram table is frozen.
